\section{Formal statements}\label{sec:lean}
\begingroup\raggedright\setlength{\parindent}{1.5em}

The main existence, continuity, and characterization results are proved in
Lean~4 (\texttt{v4.35.0}) with Mathlib in
\url{https://github.com/kejace/cd4vt} at version \texttt{\leancommit}, in the
library \texttt{Gluck}. Each displayed theorem depends only on the standard
axioms \texttt{propext}, \texttt{Classical.choice} and \texttt{Quot.sound}.
After \texttt{lake build}, the command
\texttt{lake env lean scripts/axioms.lean} prints these axioms for every
theorem below, using \texttt{\#print axioms}, for example
\begin{leancode}
#print axioms Gluck.Discrete.tabachnikov_converse_geometric
\end{leancode}
The paper sources name the Lean declarations behind the mathematical statements.

\subsection*{Theorem~\ref{thm:main}}
File \nolinkurl{Gluck/Discrete/TabachnikovConverse.lean}. Here
\texttt{vertexCircumradius v i} is the distance from \texttt{v i} to the
circumcenter of \texttt{v (i-1)}, \texttt{v i}, \texttt{v (i+1)}, and
\texttt{IsCoherentConvexPolygon} is simplicity, strict convexity and
coherence.
\begin{leancode}
theorem tabachnikov_converse_geometric [NeZero n]
    (hn : 4 ≤ n) {r : ZMod n → ℝ}
    (hr : ∀ i, 0 < r i) (hext : IsPlateauTwoMaxTwoMin r) :
    ∃ v : ZMod n → ℂ, IsCoherentConvexPolygon v ∧
      ∀ i, vertexCircumradius v i = r i
\end{leancode}

For this same polygon, the lemmas
\nolinkurl{IsCoherentConvexPolygon.signedMengerR2_eq_inv_of_radii} and
\nolinkurl{IsCoherentConvexPolygon.not_concyclic4_of_radius_ne}, in
\nolinkurl{Gluck/Discrete/TabachnikovCharacterization.lean}, give curvature
$1/r_i$ and exclude four consecutive concyclic vertices at each individual
pair with $r_i\ne r_{i+1}$. The original theorem
\texttt{tabachnikov\_converse}, with the combined conclusions, is retained.

\subsection*{Corollary~\ref{cor:iff}}
File \nolinkurl{Gluck/Discrete/TabachnikovCharacterization.lean}. The
hypotheses include positivity and nonconstancy, and the conclusion uses the
same distinct plateau runs as Theorem~\ref{thm:main}.
\begin{leancode}
theorem tabachnikov_radius_characterization [NeZero n]
    (hn : 4 ≤ n) {r : ZMod n → ℝ}
    (hr : ∀ i, 0 < r i) (hnc : ¬ ∃ c, ∀ i, r i = c) :
    (∃ v : ZMod n → ℂ, IsCoherentConvexPolygon v ∧
      ∀ i, vertexCircumradius v i = r i) ↔ IsPlateauTwoMaxTwoMin r
\end{leancode}

\subsection*{Theorem~\ref{thm:side-tangent-converse}}
File \nolinkurl{Gluck/Discrete/SideTangentConverse.lean}, namespace
\texttt{Gluck.Discrete.SideTangent}. The predicate
\texttt{SideTangentAdmissible r} means positive radii which are constant
or whose incoming-minus-outgoing differences have a strict four-sign window.
The predicate \texttt{RealizesSideTangent r} includes simplicity, strict
support, outward unit normals, all three signed-distance equations,
uniqueness of center and radius, and the open central contact.
\begin{leancode}
theorem sideTangentConverse [NeZero n] (hn : 4 ≤ n)
    (r : ZMod n → ℝ) (h : SideTangentAdmissible r) :
    RealizesSideTangent r
\end{leancode}
The same file proves \texttt{IsSideTangentRealization.tangent\_sphere}:
the sphere meets each of the three supporting lines in exactly its stated
contact point. Thus the linear equations imply metric circle tangency.

\subsection*{Corollary~\ref{cor:side-tangent-iff}}
File \nolinkurl{Gluck/Discrete/SideTangentForward.lean}, in the same
namespace and proof branch. The forward implication starts from an actual
\texttt{IsSideTangentRealization}; the balanced polygon and the necessary
sign condition are proved from it.
\begin{leancode}
theorem realizesSideTangent_iff_admissible [NeZero n]
    (hn : 4 ≤ n) (r : ZMod n → ℝ) :
    RealizesSideTangent r ↔ SideTangentAdmissible r
\end{leancode}
For positive nonconstant radii, \texttt{realizesSideTangent\_iff\_dfv}
removes the constant alternative and concludes
\texttt{DFV (radiusDifference r) 0} exactly.

\subsection*{Theorem~\ref{thm:local}}
File \nolinkurl{Gluck/Discrete/LocalContinuity.lean}. A
\texttt{LocalDiscreteRealizationFamily κ₀} is an open neighbourhood of
\texttt{κ₀} with a $C^1$ length map whose values are realizations and whose
vertex map is $C^1$. The strict window already implies nonconstancy and
$n\ge4$; neither needs to be a separate hypothesis.
\begin{leancode}
theorem localContinuousRealization_of_dfv [NeZero n]
    {κ₀ : ZMod n → ℝ} {c : ℝ}
    (hκ₀ : ∀ i, 0 < κ₀ i) (hdfv : DFV κ₀ c) :
    Nonempty (LocalDiscreteRealizationFamily κ₀)
\end{leancode}

\subsection*{Theorem~\ref{thm:hexagon}}
File \nolinkurl{Gluck/Discrete/GlobalContinuity/EnclosingCircleObstruction.lean},
where \texttt{bandProfile s} is $\kappa^s$ of~\eqref{eq:band}.
\begin{leancode}
theorem no_continuous_regular_hexagon_path :
    ¬ ∃ p : ℝ → (ZMod 6 → ℂ), ContinuousOn p (Set.Icc 0 1) ∧
      ∀ s ∈ Set.Icc (0 : ℝ) 1,
        IsSimplePolygon (p s) ∧ DahlbergRegular (p s) ∧
        SignedMengerProfile (p s) = bandProfile s
\end{leancode}
Here \texttt{DahlbergRegular} asks that at every vertex the circumcenter lie
in the closed vertex cone, with a collinear alternative when the middle
point lies between the others. The prescribed positive profile excludes
that alternative and already forces every consecutive triple to be positively
oriented. The function \texttt{SignedMengerProfile} is
$i\mapsto M(P_{i-1},P_i,P_{i+1})$. The orientation of a continuous path of
convex hexagons is constant, and complex conjugation reverses it; after
choosing counterclockwise orientation, a coherent convex realization of
$\kappa^s$ satisfies the three displayed conditions. Thus the formal statement
implies Theorem~\ref{thm:hexagon}.
\par\endgroup
